% =====================================================================
\chapter{Giới thiệu}
\label{ch:intro}
% =====================================================================

\section{Mục đích tài liệu}
Tài liệu này đặc tả yêu cầu phần mềm cho hệ thống \textbf{LabelX}, phạm vi
chức năng \textbf{M13 — Reviewer Prioritization Assistant} (tên đề tài đăng ký:
Annotation QC Studio): một trợ lý giúp reviewer quyết định
\emph{nên xem frame nào trước} dựa trên rủi ro có lỗi annotation và bằng chứng
đi kèm, thay vì review tuần tự toàn bộ dữ liệu.

Tài liệu dùng để:
\begin{itemize}
  \item thống nhất phạm vi nghiệm thu giữa Product Owner, Quality Assurance,
        đội backend và đội mô hình;
  \item làm căn cứ thiết kế, lập trình, kiểm thử và nghiệm thu;
  \item cố định \emph{cách đo} hai kết quả chính trước khi chạy pilot, để
        tránh điều chỉnh phương pháp sau khi thấy số liệu.
\end{itemize}

\section{Phạm vi sản phẩm}
\label{sec:scope}
M13 là một lát cắt (vertical slice) của module Quality Control trong LabelX.
Bản nghiệm thu đầu tiên tập trung vào \textbf{một luồng duy nhất}:

\begin{center}
\begin{tikzpicture}[node distance=4mm]
  \node[boxb, text width=22mm] (a) {Annotation\\từ CVAT};
  \node[boxb, text width=22mm, right=of a] (b) {Phân tích\\nghi vấn};
  \node[boxb, text width=22mm, right=of b] (c) {Xếp hạng\\frame};
  \node[boxb, text width=24mm, right=of c] (d) {Reviewer xem\\bằng chứng \& phân xử};
  \node[boxb, text width=22mm, right=of d] (e) {Báo cáo\\hiệu quả};
  \draw[arr] (a)--(b); \draw[arr] (b)--(c); \draw[arr] (c)--(d); \draw[arr] (d)--(e);
\end{tikzpicture}
\end{center}

\begin{xltabular}{\linewidth}{@{}L{3.1cm}Y@{}}
\caption{Phân loại phạm vi của bản nghiệm thu M13}\label{tab:scope}\\
\toprule
\textbf{Nhóm} & \textbf{Nội dung}\\
\midrule
\endfirsthead
\toprule \textbf{Nhóm} & \textbf{Nội dung}\\ \midrule
\endhead
\bottomrule
\endfoot
\textbf{Lõi — nghiệm thu} &
Đọc annotation và ảnh từ CVAT, tạo snapshot bất biến; chạy các engine phát hiện
nghi vấn cho Bounding box; hợp nhất candidate; tính điểm rủi ro frame và xếp
hạng; hàng đợi ưu tiên; Review Workspace hiển thị bằng chứng; quyết định
xác nhận/bác bỏ/chưa chắc chắn; phân xử; đo Recall@20\% và mức giảm công sức
review; báo cáo hiệu quả.\\
\textbf{Hỗ trợ — mức cần thiết} &
Tra cứu guideline theo rule ID/version; yêu cầu sửa (rework) và kiểm lại sau
sửa; phân quyền và nhật ký kiểm toán; Quality Gate tối thiểu để không phát hành
dữ liệu chưa xác minh.\\
\textbf{Ngoài phạm vi bản đầu} &
Temporal/track, polygon, polyline, video (B-09); ghi annotation ngược vào CVAT
(B-18); Diff/pre-label, vehicle profile, feedback loop (B-16); Classifier thứ
hai (B-07); Guideline RAG (B-13); phát hành dataset đầy đủ và quản trị nền tảng
ngoài mức tối thiểu.\\
\textbf{Đề xuất thu hẹp — cần duyệt} &
Engine mô hình thị giác – ngôn ngữ (VLM) không bật trong pilot M13. Đây là
\emph{đề xuất thu hẹp cho pilot M13, cần Product Owner duyệt}; nó không kế thừa
từ B-08. Quyết định B-08 (pipeline candidate $\rightarrow$ chọn theo policy
$\rightarrow$ worker kiểm chọn lọc) vẫn là kiến trúc đã chốt. Nếu đề xuất bị bác,
VLM chạy theo B-08 với hạn mức 400 candidate/run, tối đa 3 lần thử, và hết hạn
mức thì ghi phần chưa kiểm (\TBD[-08]).\\
\end{xltabular}

\section{Đối tượng đọc và cách đọc}
\begin{xltabular}{\linewidth}{@{}L{4cm}Y@{}}
\toprule
\textbf{Người đọc} & \textbf{Nên đọc}\\
\midrule
Product Owner, ban nghiệm thu & Chương~\ref{ch:overview}, \ref{ch:eval}, \ref{ch:trace}.\\
Quality Assurance Lead & Chương~\ref{ch:data}, \ref{ch:usecase}, \ref{ch:eval}.\\
Backend / Frontend & Chương~\ref{ch:usecase}–\ref{ch:arch}.\\
Đội mô hình & Chương~\ref{ch:data}, mục~\ref{sec:fr-eng}, \ref{sec:fr-rnk}, chương~\ref{ch:eval}.\\
Kiểm thử (QA engineer) & Chương~\ref{ch:func}, \ref{ch:eval}, \ref{ch:trace}.\\
\bottomrule
\end{xltabular}

\section{Quy ước tài liệu}
\begin{itemize}
  \item \textbf{Mã định danh:} \req{UC-xx} use case; \req{FR-<MODULE>-xx} yêu cầu
        chức năng; \req{NFR-xx} yêu cầu phi chức năng; \req{AC-xx} tiêu chí
        nghiệm thu; \req{BR-xx} quy tắc nghiệp vụ; \req{B-xx} quyết định kiến
        trúc đã chốt; \req{R-xx} mục tiêu kiểm chứng đã chốt; \req{E1–E3} nhóm lỗi.
  \item \textbf{Mức ưu tiên (MoSCoW):} \must{} = Must (bắt buộc cho nghiệm thu),
        \should{} = Should (nên có), \could{} = Could (có thể hoãn).
  \item \textbf{Từ khoá:} “\emph{phải}” = bắt buộc; “\emph{nên}” = khuyến nghị;
        “\emph{có thể}” = tuỳ chọn.
  \item \TBD{} đánh dấu tham số chưa chốt. Mỗi \TBD{} có người chịu trách
        nhiệm trong bảng~\ref{tab:tbd}.
  \item Thuật ngữ kỹ thuật giữ nguyên tiếng Anh khi đó là tên dùng trong sản
        phẩm (ví dụ \term{snapshot}, \term{candidate}); giải thích ở mục~\ref{sec:glossary}.
\end{itemize}

\section{Tài liệu tham chiếu}
\begin{enumerate}[label={[\arabic*]}, leftmargin=2.2em]
  \item ISO/IEC/IEEE 29148:2018 — \emph{Systems and software engineering — Life cycle processes — Requirements engineering}.
  \item F.~Yu và cộng sự, \emph{BDD100K: A Diverse Driving Dataset for Heterogeneous Multitask Learning}, CVPR 2020.
  \item CVAT — Computer Vision Annotation Tool, tài liệu REST API (phiên bản cụ thể \TBD[-01]).
  \item LabelX Design System và 25 màn hình (\code{docs/design/}) — nguồn H.
  \item Bản rà soát kiến trúc, luồng và triển khai đã chốt, 05/10/2026 — nguồn R.
  \item QC Engine \& Review System — đề xuất kiến trúc, 04/10/2026 — nguồn T.
  \item Quality Control Review — UX Architecture — nguồn A.
\end{enumerate}

\section{Thuật ngữ và viết tắt}
\label{sec:glossary}

\begin{xltabular}{\linewidth}{@{}L{3.6cm}Y@{}}
\caption{Bảng thuật ngữ}\label{tab:glossary}\\
\toprule
\textbf{Thuật ngữ} & \textbf{Định nghĩa trong tài liệu này}\\
\midrule
\endfirsthead
\toprule \textbf{Thuật ngữ} & \textbf{Định nghĩa}\\ \midrule
\endhead
\midrule \multicolumn{2}{r}{\footnotesize\itshape tiếp trang sau}\\
\endfoot
\bottomrule
\endlastfoot
\multicolumn{2}{@{}l}{\textbf{\sffamily Dữ liệu và annotation}}\\
Frame & Một ảnh là đơn vị review và xếp hạng. Với BDD100K detection, mỗi ảnh tĩnh là một frame.\\
Annotation & Một đối tượng được gán nhãn trên frame: Bounding box + lớp + thuộc tính.\\
Bounding box (BBox) & Hình chữ nhật song song trục \((x_1, y_1, x_2, y_2)\) theo pixel.\\
Taxonomy & Danh sách lớp và thuộc tính hợp lệ của dataset, có version.\\
Guideline & Hướng dẫn gán nhãn, chia thành rule có mã (ví dụ \code{VEH-03}) và version.\\
CVAT & Công cụ gán nhãn nguồn. Annotation chỉ được sửa trên CVAT (B-18).\\
Project / Task / Job & Cấp tổ chức dữ liệu trong CVAT; Job là phần việc giao cho một annotator.\\
Revision & Trạng thái annotation tại một thời điểm, nhận diện bằng hash nội dung đã chuẩn hoá.\\
Snapshot & Bản sao bất biến của annotation, metadata và ảnh tại một revision; mọi phân tích chạy trên snapshot (B-02).\\
Drift & Annotation trên CVAT thay đổi trong lúc đang export snapshot; phải phát hiện và chặn khoá.\\
\multicolumn{2}{@{}l}{\textbf{\sffamily Phân tích và xếp hạng}}\\
Engine & Thành phần kiểm tra tạo candidate. M13 dùng Schema/Taxonomy, Geometry, Duplicate/Overlap, Mô hình độc lập (Detector) và Metric.\\
QC Run & Một lần chạy các engine trên một snapshot với cấu hình, seed và version cố định.\\
Detector & Mô hình phát hiện đối tượng độc lập với người gán nhãn, dùng sinh candidate thiếu box và sai lớp. Không phải Ground Truth.\\
Matching & Ghép một-một giữa box của Detector (hoặc reference) với annotation theo IoU, hình học trước rồi mới so lớp (B-21).\\
IoU & Intersection over Union — tỉ lệ diện tích giao trên diện tích hợp của hai box.\\
Candidate & Nghi vấn do engine sinh ra, chưa được người xác nhận. Candidate không phải lỗi.\\
Issue & Đơn vị công việc review sau khi gộp các candidate trùng (B-10). Có vòng đời riêng.\\
Evidence (bằng chứng) & Dữ liệu giải thích candidate: box của Detector, confidence, IoU, rule vi phạm, guideline, case tương tự.\\
Risk score \(s(f)\) & Điểm rủi ro của frame \(f\), dùng để xếp hạng. Có version công thức.\\
Ranking \(\pi\) & Thứ tự frame giảm dần theo \(s(f)\), có quy tắc phá hoà cố định.\\
Top-\(k\)\% & Tập \(\lceil k\cdot N/100\rceil\) frame đứng đầu ranking trong tập \(N\) frame.\\
Coverage ledger & Sổ ghi phần đơn vị đã kiểm/chưa kiểm theo từng engine với mẫu số áp dụng (B-04).\\
Not checked / Partial / Failed & Trạng thái công khai của engine: chưa kiểm, kiểm một phần, lỗi thực thi. Không bao giờ được coi là đạt (B-19).\\
\multicolumn{2}{@{}l}{\textbf{\sffamily Review và quy trình}}\\
Review queue & Hàng đợi issue/frame cho reviewer, tách theo nguồn (rủi ro, ngẫu nhiên).\\
Review Workspace & Màn hình xem frame, annotation, bằng chứng và ra quyết định.\\
Quyết định review & Một trong: Xác nhận lỗi, Bác bỏ, Chưa chắc chắn, Chuyển cấp trên.\\
False alarm (cảnh báo sai) & Candidate/issue bị reviewer bác bỏ hoặc phân xử là không lỗi.\\
Escalation / Adjudication & Chuyển issue chưa đủ căn cứ cho Quality Assurance Lead phân xử.\\
Guideline Gap & Ghi nhận chỗ guideline thiếu, phát sinh từ phân xử.\\
Rework & Yêu cầu annotator sửa trên CVAT; tạo revision mới.\\
Re-check / Verify & Chạy lại engine bị ảnh hưởng trên revision mới và reviewer xác minh trước khi đóng issue.\\
QC run cuối & QC Run trên snapshot của revision đã sửa; mọi báo cáo và phát hành trỏ về run này (B-03).\\
Lease & Quyền giữ tạm một issue/frame cho một reviewer, tránh hai người cùng xử lý.\\
Self-review & Reviewer review annotation do chính mình gán; backend phải chặn (B-12).\\
Waiver & Miễn trừ có thời hạn cho một điều kiện gate, cần người duyệt khác người yêu cầu (B-11).\\
Quality Gate & Tập điều kiện phải đạt trước khi một revision được coi là dùng được.\\
\multicolumn{2}{@{}l}{\textbf{\sffamily Đánh giá}}\\
Reference & Gồm annotation chuẩn (GT) của toàn bộ tập đánh giá và tập lỗi suy ra từ GT; cùng khoá version trước khi đo (B-20, mục~\ref{sec:reference}).\\
Ground Truth & Annotation chuẩn đã được chuyên gia duyệt và khoá. Không model nào được mặc định là Ground Truth.\\
Lỗi đã xác minh & Một phần tử của tập lỗi \(\mathcal{E}\): bản ghi gắn một nhóm lỗi với một đối tượng GT (E1, E2) hoặc một annotation dư (E3), suy ra từ GT và được người xác minh duyệt.\\
Tập hiệu chỉnh / held-out & Tập dùng tinh chỉnh công thức điểm và tập chỉ dùng đo nghiệm thu; không giao nhau.\\
Recall@20\% & Tỉ lệ lỗi đã xác minh nằm trong 20\% frame đầu của ranking (mục~\ref{sec:kpi1}).\\
Effort & Tổng thời gian thao tác có ích của reviewer/QA, gồm xử lý cảnh báo sai và kiểm lại sau sửa (mục~\ref{sec:kpi2}).\\
Baseline & Quy trình review thủ công hiện tại, không có ranking và bằng chứng.\\
Assisted & Quy trình review có trợ lý M13.\\
Residual error rate & Tỉ lệ lỗi còn lại sau quy trình, đo trên reference.\\
Non-inferiority & Kiểm định chất lượng nhánh assisted không kém baseline quá biên \(\delta\).\\
Bootstrap & Phương pháp lấy mẫu lại để ước lượng khoảng tin cậy.\\
Model Orchestrator & Thành phần tổng hợp kết quả đánh giá engine trên reference, kèm provenance; không dùng model làm trọng tài (B-05). Đặc tả chi tiết \TBD[-21].\\
\multicolumn{2}{@{}l}{\textbf{\sffamily Viết tắt}}\\
QA / QC & Quality Assurance / Quality Control.\\
SRS & Software Requirements Specification.\\
RBAC & Role-Based Access Control.\\
VLM & Vision–Language Model — mô hình thị giác – ngôn ngữ.\\
DRF & Django REST framework.\\
MVP & Minimum Viable Product.\\
KPI & Key Performance Indicator.\\
\end{xltabular}
